% ============================================================
% РАЗДЕЛ 1. Проблема холодного старта ML-моделей
% ============================================================
\section{ПРОБЛЕМА ХОЛОДНОГО СТАРТА ML-МОДЕЛЕЙ В БЕССЕРВЕРНОМ ИНФЕРЕНСЕ}

Настоящий раздел посвящён обоснованию актуальности проблемы холодного старта в системах бессерверного инференса моделей машинного обучения. Рассматривается эволюция архитектур и размеров ML-моделей, показывающая, что объём весовых файлов вырос от единиц мегабайт до сотен гигабайт. Раскрывается механизм каждого этапа холодного старта в Kubernetes, обосновывается, почему классическая схема с init-контейнером и томом \texttt{emptyDir} неэффективно использует уже загруженные данные. Оцениваются последствия для соблюдения SLA и пользовательского опыта с применением закона Литтла и вероятностной модели холодного старта. Формулируются функциональные и нефункциональные требования к системе оптимизации.

% ============================================================
\subsection{Эволюция размеров моделей машинного обучения}
% ============================================================

\subsubsection{От компактных сетей к масштабным языковым моделям}

На раннем этапе развития глубокого обучения доминировали сверточные нейронные сети (CNN), ориентированные преимущественно на задачи компьютерного зрения. Прорывная работа Крижевского, Суцкевера и Хинтона представила модель AlexNet в 2012~году: сеть насчитывала 60~миллионов параметров, а объем её весовых файлов в 32-битном формате составлял около 233~МБ~\cite{alexnet2012}. При типичной для того времени пропускной способности сети порядка 100~Мбит/с такой объём передавался за несколько секунд и не создавал значимых инфраструктурных проблем. Параллельно развивались рекуррентные нейронные сети (RNN) и их разновидности LSTM и GRU для задач обработки последовательностей, однако их размеры также оставались в пределах сотен мегабайт.

Качественный переход произошел в 2017~году с публикацией архитектуры Transformer. Механизм самовнимания (self-attention) обеспечил возможность эффективно обрабатывать длинные контекстные зависимости и параллелизовать обучение, что сделало Transformer доминирующей архитектурой для задач обработки естественного языка~\cite{bert2018}. Вместе с этим значительно выросли и вычислительные требования.

Модель BERT~Base (2018) с 110~миллионами параметров занимала около 440~МБ в 32-битном формате; версия BERT~Large содержала 340~миллионов параметров и требовала порядка 1.4~ГБ~\cite{bert2018}. Уже на этом этапе загрузка весов из удалённого хранилища начинала ощутимо влиять на задержку инициализации сервиса.

\subsubsection{Экспоненциальный рост весов языковых моделей}

Выход GPT-3 в 2020~году ознаменовал качественный скачок: 175~миллиардов параметров потребовали около 350~ГБ дискового пространства для хранения весов в 16-битном формате~\cite{gpt3_2020}. Для загрузки такого объёма по каналу 10~Гбит/с требуется около 280~секунд — почти пять минут только на передачу данных. В реальных условиях с конкурентными запросами из других подов это время возрастает существенно.

Тенденция к открытым моделям, усиленная развитием семейства DeepSeek, расширила доступ исследователей и производственных команд к крупным моделям. В работе DeepSeek-LLM рассматриваются конфигурации 7B и 67B параметров, а расчётный объём основных весов варианта DeepSeek-LLM~67B в формате FP16 составляет около 134~ГБ~\cite{deepseekllm_2024}. Следующее поколение, DeepSeek~V2, уже переходит в диапазон сотен миллиардов параметров~\cite{deepseekv2_2024}.

\begin{table}[htbp]
\centering
\caption{Расчётный объём весов ключевых ML-моделей в формате FP16}
\label{tab:model_sizes}
\scriptsize
\begin{tabularx}{\textwidth}{|>{\raggedright\arraybackslash}p{24mm}|>{\centering\arraybackslash}p{10mm}|>{\raggedright\arraybackslash}p{23mm}|>{\raggedleft\arraybackslash}p{18mm}|>{\raggedleft\arraybackslash}p{18mm}|>{\raggedright\arraybackslash}X|}
\hline
\textbf{Модель} & \textbf{Год} & \textbf{Архитектура} & \textbf{Параметры} & \textbf{Вес (FP16)} & \textbf{Влияние} \\ \hline
AlexNet          & 2012 & CNN         & 60~млн    & $\approx$0.12~ГБ & Незначительное \\
\hline
BERT Base        & 2018 & Transformer & 110~млн   & $\approx$0.22~ГБ & Малое \\
\hline
GPT-2 XL         & 2019 & Transformer & 1.5~млрд  & $\approx$3~ГБ    & Умеренное \\
\hline
T5-11B           & 2020 & Transformer & 11~млрд   & $\approx$22~ГБ   & Существенное \\
\hline
GPT-3            & 2020 & Transformer & 175~млрд  & $\approx$350~ГБ  & Критическое \\
\hline
DeepSeek-LLM~7B  & 2024 & Transformer & 7~млрд    & $\approx$14~ГБ   & Существенное \\
\hline
DeepSeek-LLM~67B & 2024 & Transformer & 67~млрд   & $\approx$134~ГБ  & Критическое \\
\hline
Falcon~180B      & 2023 & Transformer & 180~млрд  & $\approx$360~ГБ  & Критическое \\
\hline
DeepSeek~V2     & 2024 & MoE         & 235.7~млрд & $\approx$472~ГБ & Критическое \\
\hline
GLM-4.5         & 2025 & MoE         & 358.3~млрд & $\approx$717~ГБ & Критическое \\
\hline
\shortstack[l]{DeepSeek\\V4 Pro} & 2026 & MoE & 861.6~млрд & $\approx$1.72~ТБ & Критическое \\
\hline
\end{tabularx}
\end{table}

Значения в таблице получены как произведение числа параметров на два байта, то есть описывают только основные тензоры в FP16 и используют десятичный гигабайт ($10^9$ байт). Фактический размер репозитория зависит от формата сериализации, разбиения на файлы и служебных данных. Для Mixtral~8$\times$22B приведены все 141~млрд параметров, хотя на одном токене активируется 39~млрд. Число параметров приведено по первичным публикациям и карточкам моделей~\cite{bert2018,gpt2_2019,t5_2020,gpt3_2020,deepseekllm_2024,falcon2023,mixtral2024,deepseekv2_2024,glm45_2025,deepseekv4_2026}.

Таблица~\ref{tab:model_sizes} демонстрирует, что за период с 2018 по 2026~год расчётный объём весов ведущих моделей вырос от сотен мегабайт до терабайт --- почти на четыре порядка. Если для BERT~Base загрузка занимала секунды, то для DeepSeek~V4~Pro при той же пропускной способности сети она может занимать часы. Именно это изменение масштаба превращает загрузку весов из тривиальной операции в доминирующий компонент задержки холодного старта.

\begin{figure}[htbp]
\centering
\includegraphics[width=0.88\textwidth]{chart_model_sizes.png}
\caption{Монотонная последовательность отобранных ML-моделей с возрастающим расчётным объёмом основных весов в FP16, 2018--2026 (логарифмическая шкала). График иллюстрирует рост выбранной последовательности, а не утверждает, что каждый выпуск соответствующего года крупнее всех других моделей.}
\label{fig:model_sizes}
\end{figure}

\subsubsection{Законы масштабирования как движущая сила роста}

Рост числа параметров подчиняется так называемым законам масштабирования (scaling laws). Эмпирически установлено, что качество языковых моделей, измеряемое кросс-энтропийными потерями, улучшается в степенной зависимости от трёх факторов: объёма параметров модели, размера обучающего набора данных и доступных вычислительных ресурсов~\cite{gpt3_2020}. Из этих законов следует устойчивый технологический стимул к дальнейшему увеличению масштаба: большие модели достигают того же уровня качества с меньшим числом шагов обучения и могут демонстрировать способности, не наблюдаемые у меньших моделей.

Параллельно с ростом параметров развивается концепция «плотности возможностей» (capability density)~\cite{capability_density}: компактные модели, такие как Mistral~7B и семейство Phi-3, достигают результатов, сопоставимых с более крупными аналогами, благодаря качественным обучающим данным и архитектурным оптимизациям. Однако даже модели в диапазоне 7--70~млрд параметров требуют примерно от 13 до 140~ГБ для хранения основных весов в FP16 --- объёмов, принципиально меняющих характер инфраструктурной задачи.

Таким образом, рост размеров ML-моделей является устойчивой долгосрочной тенденцией, а значит, проблема загрузки весов при холодном старте будет лишь обостряться с течением времени.

% ============================================================
\subsection{Механизм холодного старта в Kubernetes}
% ============================================================

\subsubsection{Декомпозиция задержки холодного старта}

При развёртывании ML-инференса в Kubernetes каждый новый инстанс сервиса проходит ряд последовательных фаз. Суммарная задержка холодного старта формализуется следующим образом:

\begin{equation}
T_{\mathrm{start}} = T_{\mathrm{schedule}} + T_{\mathrm{pull}} + T_{\mathrm{unpack}} + T_{\mathrm{init}} + T_{\mathrm{download}} + T_{\mathrm{load}},
\label{eq:cold_start}
\end{equation}

где $T_{\mathrm{schedule}}$ --- время планировщика Kubernetes для назначения пода на конкретную ноду; $T_{\mathrm{pull}}$ --- загрузка образа контейнера из реестра; $T_{\mathrm{unpack}}$ --- декомпрессия и распаковка слоёв образа на диск; $T_{\mathrm{init}}$ --- выполнение init-контейнера (в том числе storage-initializer); $T_{\mathrm{download}}$ --- скачивание весов модели из внешнего репозитория; $T_{\mathrm{load}}$ --- загрузка весов из файловой системы в оперативную память и GPU-видеопамять.

Каждое слагаемое вносит самостоятельный вклад в суммарную задержку. Рассмотрим их характеристики и причины, по которым традиционная схема не позволяет эффективно сократить наиболее затратные фазы.

\subsubsection{Фаза планирования: $T_{\mathrm{schedule}}$}

Планировщик Kubernetes анализирует доступные ноды с учётом запрошенных ресурсов (CPU, память, GPU), ограничений через nodeSelector и tolerations, правил affinity и topologySpreadConstraints. При загруженном кластере или ожидании освобождения GPU-слота эта фаза может занимать от долей секунды до нескольких минут. Поскольку планировщик не учитывает наличие кэшированных весов модели на ноде, он может назначить под на ноду без нужных данных, вынуждая к их повторной загрузке.

\subsubsection{Фазы загрузки и распаковки образа: $T_{\mathrm{pull}} + T_{\mathrm{unpack}}$}

Образ контейнера для ML-инференса включает среду выполнения языка, фреймворк (PyTorch, TensorFlow), CUDA-библиотеки и вспомогательные зависимости. Для типичного контейнера с vLLM или Triton Inference Server размер образа может составлять 15--25~ГБ. Загрузка по каналу 10~Гбит/с занимает ориентировочно 12--20~секунд для передачи данных, однако реальное время больше из-за последовательного характера протокола реестра контейнеров: слои образа обычно скачиваются без полноценного параллелизма~\cite{bentoml2024cold}.

После скачивания каждый слой образа декомпрессируется из формата gzip или zstd и записывается на диск. Эта операция CPU-интенсивна: декомпрессия 20~ГБ образа может занять от трёх до восьми минут. Механизм Copy-on-Write, используемый большинством драйверов хранения контейнеров (overlay2, fuse-overlayfs), добавляет накладные расходы при последующем доступе к файлам внутри образа.

\subsubsection{Фаза init-контейнера: $T_{\mathrm{init}}$}

Классическая схема развёртывания ML-инференса в Kubernetes предполагает использование init-контейнера (storage-initializer), который запускается до основного runtime-контейнера. Его задача --- скачать файлы весов модели из внешнего хранилища (HuggingFace Hub, S3, GCS) и записать их в том типа \texttt{emptyDir}, разделяемый с основным контейнером.

Принципиальная схема взаимодействия выглядит следующим образом: планировщик назначает под на ноду, после чего инициализируется том \texttt{emptyDir} на локальном диске этой ноды. Затем запускается init-контейнер \texttt{storage-initializer}, который последовательно скачивает файлы весов и записывает их в смонтированный \texttt{emptyDir}. После успешного завершения init-контейнера запускается runtime-контейнер (vLLM, TGI, Triton), который считывает файлы из \texttt{emptyDir} и загружает их в GPU-память.

\subsubsection{Фаза загрузки весов из репозитория: $T_{\mathrm{download}}$}

Компонент $T_{\mathrm{download}}$ является наиболее критичным с точки зрения зависимости от размера модели. При линейном росте объёма весов это слагаемое растёт пропорционально. Для DeepSeek-LLM~7B (порядка 15~ГБ) при канале 1~Гбит/с без конкуренции загрузка занимает ориентировочно 120~секунд; для DeepSeek-LLM~67B (порядка 140~ГБ) --- уже порядка 1120~секунд, то есть около 19~минут~\cite{bentoml2024cold}.

Ключевая особенность проблемы состоит в том, что при масштабировании до $N$ реплик на разных нодах внешний репозиторий получает $N$ независимых запросов на один и тот же набор файлов. Каждый init-контейнер выполняет полную загрузку, игнорируя тот факт, что те же данные могут уже присутствовать на соседних нодах кластера.

\subsubsection{Ограничения схемы с \texttt{emptyDir}}

Том \texttt{emptyDir} привязан к жизненному циклу пода и физически размещается на диске конкретной ноды. При завершении или перепланировании пода данные в \texttt{emptyDir} безвозвратно удаляются. Следствием является отсутствие механизма разделения данных между репликами на одной ноде: при масштабировании к двум репликам на одном узле каждая из них независимо скачивает полный набор весов. Отсутствует и возможность передачи уже загруженных данных между нодами кластера, что исключает локальное повторное использование при перепланировании подов.

Важно разграничить два типа повторной загрузки. Первый --- внутринодовая избыточность: несколько подов на одной ноде, работающих с одной и той же моделью, независимо скачивают и хранят идентичные веса на одном диске. Второй --- межнодовая избыточность: если модель уже загружена на соседней ноде, схема с \texttt{emptyDir} не предоставляет никакого механизма для получения этих данных без повторного обращения во внешний репозиторий. При масштабировании системы до $N$ реплик, распределённых по $M$ нодам, возникает до $N$ независимых операций загрузки одного и того же содержимого, что создаёт кратную нагрузку на канал до внешнего репозитория и увеличивает операционные расходы.

\subsubsection{Фаза загрузки в GPU-память: $T_{\mathrm{load}}$}

После того как файлы весов оказались в \texttt{emptyDir}, runtime-контейнер (vLLM, TGI, Triton) считывает их с диска и загружает в GPU-память. Для модели класса DeepSeek-LLM~7B эта операция занимает ориентировочно 60--120~секунд; для крупных моделей --- несколько минут. Кроме того, некоторые inference-движки, в частности vLLM, выполняют компиляцию графа CUDA и профилирование KV-кэша, что дополнительно увеличивает $T_{\mathrm{load}}$~\cite{bentoml2024cold}.

\subsubsection{Итоговый характер задержки}

Для типичного контейнера с моделью семейства DeepSeek объёмом порядка 15~ГБ суммарная задержка холодного старта на чистом узле складывается ориентировочно следующим образом: $T_{\mathrm{pull}} + T_{\mathrm{unpack}} \approx 5$--8~мин; $T_{\mathrm{download}} \approx 2$--5~мин; $T_{\mathrm{load}} \approx 1$--3~мин. Итоговое время в зависимости от конфигурации сети и дисков может превышать 10--15~минут~\cite{bentoml2024cold}.

При этом задержка масштабируется нелинейно: при переходе от 8B к 70B параметров $T_{\mathrm{download}}$ возрастает пропорционально размеру весов, тогда как $T_{\mathrm{load}}$ дополнительно зависит от скорости шины памяти GPU и конкретного inference-движка.

\begin{figure}[htbp]
\centering
\includegraphics[width=0.88\textwidth]{chart_cold_start_breakdown.png}
\caption{Декомпозиция задержки холодного старта для DeepSeek-LLM~7B (15~ГБ, канал 200~МБ/с). Фаза загрузки весов ($T_{\mathrm{download}} \approx 75$~с) составляет около 80\% суммарного времени~$\approx 94$~с.}
\label{fig:cold_start_breakdown}
\end{figure}

% ============================================================
\subsection{Влияние холодного старта на SLA, пользовательский опыт и стоимость}
% ============================================================

\subsubsection{SLA и метрики задержки}

Соглашение об уровне обслуживания (SLA) для интерактивных AI-сервисов, как правило, регламентирует задержку на перцентилях P95 и P99, а также специфические для генеративных моделей метрики: TTFT (Time to First Token --- время до первого токена) и TPS (Tokens Per Second --- скорость генерации). Для диалоговых ассистентов требования к TTFT на P95 нередко составляют от нескольких сотен миллисекунд до нескольких секунд, а для голосовых систем порог ещё ниже.

Холодный старт, занимающий 10--15~минут, делает соблюдение таких требований при использовании стратегии scale-to-zero принципиально невозможным без специальных архитектурных мер. Запросы, поступившие в момент инициализации пода, вынуждены ожидать завершения всей цепочки фаз из формулы~(\ref{eq:cold_start}), что многократно нарушает целевые значения P95 и P99 latency.

\subsubsection{Закон Литтла: накопление ожидающих запросов}

Для количественной оценки влияния холодного старта применим закон Литтла~\cite{little1961}. Рассмотрим подсистему ожидания инициализации подов как стационарную систему массового обслуживания. Если $L_{\mathrm{wait}}$ --- среднее число запросов, находящихся в состоянии ожидания завершения холодного старта, $\lambda$ --- интенсивность входящих запросов в период масштабирования, а $W_{\mathrm{cold}}$ --- среднее время холодного старта, то:

\begin{equation}
L_{\mathrm{wait}} = \lambda \cdot W_{\mathrm{cold}}.
\label{eq:little}
\end{equation}

Из формулы~(\ref{eq:little}) следует фундаментальный вывод: при $\lambda = 10$~запр./с и $W_{\mathrm{cold}} = 600$~с (10~мин) в очереди единовременно накапливается $L_{\mathrm{wait}} = 6000$ ожидающих запросов. Каждый из них занимает соединение и ресурсы API-шлюза, что при достижении порогов ограничения по соединениям приводит к отказу в обслуживании.

Закон Литтла носит универсальный характер: он не зависит от распределения времён поступления запросов и от внутренней структуры системы массового обслуживания~\cite{little1961}. Из уравнения~(\ref{eq:little}) прямо следует, что уменьшение $W_{\mathrm{cold}}$ в $k$ раз пропорционально уменьшает число ожидающих запросов $L_{\mathrm{wait}}$. Это делает сокращение времени холодного старта первоочередной задачей при проектировании масштабируемого ML-инференса.

\subsubsection{Вероятность холодного старта при стратегии scale-to-zero}

В бессерверных системах с автомасштабированием до нуля применяется механизм keep-alive: реплика сервиса поддерживается в активном состоянии в течение времени $T_{\mathrm{keep}}$ после последнего запроса, после чего освобождается. Если поток входящих запросов к конкретной модели описывается пуассоновским процессом с интенсивностью $\lambda$, то вероятность холодного старта для произвольного запроса определяется формулой~\cite{little1961}:

\begin{equation}
P_{\mathrm{cold}} = e^{-\lambda \cdot T_{\mathrm{keep}}},
\label{eq:pcold}
\end{equation}

где $T_{\mathrm{keep}}$ --- время поддержания реплики в активном состоянии после последнего запроса.

При малой нагрузке $\lambda \to 0$, и $P_{\mathrm{cold}} \to 1$: практически каждый запрос инициирует холодный старт. Это характерно для длинного «хвоста» редко используемых специализированных моделей. Обратно, при высокой нагрузке $\lambda \to \infty$, и $P_{\mathrm{cold}} \to 0$: реплики постоянно активны и к ним поступают запросы до истечения keep-alive.

Из формулы~(\ref{eq:pcold}) следует, что при фиксированном целевом SLA на максимально допустимую вероятность холодного старта $P_{\mathrm{cold}} \leq \varepsilon$ требуемое время поддержания реплики определяется как:

\begin{equation}
T_{\mathrm{keep}} \geq -\frac{\ln \varepsilon}{\lambda}.
\label{eq:keep_alive}
\end{equation}

Для модели с $\lambda = 0.01$~запр./с и допустимым $\varepsilon = 0.05$ (5\% вероятность холодного старта) получаем $T_{\mathrm{keep}} \geq 299$~с, то есть около 5~минут. За это время простаивающая GPU потребляет энергию и занимает слот, за который взимается плата --- при типичной стоимости GPU-инстанса это превращается в значительную статью операционных расходов.

\subsubsection{Дилемма scale-to-zero и over-provisioning}

Из приведённого анализа вытекает фундаментальная дилемма, с которой сталкиваются операторы ML-платформ. Стратегия scale-to-zero обеспечивает оптимальное использование ресурсов в периоды низкой нагрузки, однако при высоком $W_{\mathrm{cold}}$ делает невозможным соблюдение SLA на задержку первого токена.

Альтернатива --- over-provisioning, постоянное поддержание пула «тёплых» реплик --- устраняет задержку, но аннулирует экономическое преимущество бессерверной модели. Для модели, которая обрабатывает в среднем один запрос в час, постоянное удержание одной активной реплики с GPU-ускорителем создаёт крайне низкую утилизацию оборудования и высокую удельную стоимость обработки запроса.

Оптимальным является третий путь: снижение $W_{\mathrm{cold}}$ до уровня, при котором допустимо использовать небольшой keep-alive или даже полный scale-to-zero без критических нарушений SLA. Если холодный старт занимает 30~секунд вместо 10~минут, то допустимое время ожидания первого запроса переходит в категорию приемлемых задержек для многих производственных сценариев.

Таблица~\ref{tab:cold_start_components} иллюстрирует ориентировочный вклад различных компонентов в суммарную задержку холодного старта для модели с 8 и 70~миллиардами параметров при типичных конфигурациях. Следует подчеркнуть, что приведённые значения носят оценочный характер и существенно зависят от пропускной способности сети, производительности дискового ввода-вывода и конкретного inference-движка.

\begin{table}[htbp]
\centering
\caption{Ориентировочный вклад фаз холодного старта (оценочные данные)}
\label{tab:cold_start_components}
\begin{tabularx}{\textwidth}{|>{\raggedright\arraybackslash}X|>{\centering\arraybackslash}p{23mm}|>{\centering\arraybackslash}p{23mm}|>{\raggedright\arraybackslash}p{28mm}|}
\hline
\textbf{Фаза} & \multicolumn{2}{c|}{\textbf{Типичная длительность}} & \textbf{Зависимость} \\
              & \textbf{8B модель} & \textbf{70B модель} & \textbf{от размера модели} \\ \hline
$T_{\mathrm{schedule}}$ (планирование) & 5--30~с   & 5--30~с   & Нет \\
$T_{\mathrm{pull}}$ (загрузка образа)  & 2--5~мин  & 2--5~мин  & Слабая \\
$T_{\mathrm{unpack}}$ (распаковка)     & 2--4~мин  & 2--4~мин  & Слабая \\
$T_{\mathrm{download}}$ (загрузка весов) & 2--4~мин & 15--25~мин & Линейная \\
$T_{\mathrm{load}}$ (в GPU-память)     & 1--2~мин  & 5--10~мин & Линейная \\ \hline
$T_{\mathrm{start}}$ (итого)           & 7--15~мин & 29--44~мин & Значимая \\
\hline
\end{tabularx}
\end{table}

Из таблицы~\ref{tab:cold_start_components} видно, что при переходе от 8B к 70B параметров компоненты $T_{\mathrm{download}}$ и $T_{\mathrm{load}}$ вносят наибольший вклад в прирост суммарной задержки. Именно на эти компоненты должны быть направлены основные архитектурные оптимизации.

\subsubsection{Влияние на пользовательский опыт}

С точки зрения восприятия пользователем задержка свыше 3~секунд при обращении к веб-сервису существенно снижает вероятность продолжения взаимодействия. Для голосовых ассистентов и интерактивных чат-ботов практически любая задержка первого токена сверх 2--3~секунд воспринимается негативно. Холодный старт ML-сервиса, длящийся несколько минут, делает продукт практически непригодным к использованию в периоды низкой нагрузки при стратегии scale-to-zero. Это вынуждает продуктовые команды выбирать между ухудшением UX и значительными дополнительными затратами на инфраструктуру.

% ============================================================
\subsection{Требования к системе оптимизации холодного старта}
% ============================================================

На основании проведённого анализа механизма холодного старта, его количественных характеристик и экономических последствий сформулированы следующие требования к системе оптимизации.

\subsubsection{Функциональные требования}

\textbf{ФТ-1. Минимизация задержки первого чтения данных модели.} Система должна предоставлять первые байты данных модели в наиболее короткие сроки, не дожидаясь полной загрузки модели из внешнего источника. Для этого необходима поддержка потоковой или ленивой загрузки с возможностью начать инференс до полного завершения скачивания.

\textbf{ФТ-2. Повторное использование уже загруженных весов.} Если файлы модели или их фрагменты уже присутствуют на дисках узлов кластера, новые реплики должны получать данные из этого локального или соседнего источника без обращения к внешнему репозиторию. Повторная загрузка идентичных байт из интернета при наличии их в кластере недопустима.

\textbf{ФТ-3. Прозрачность для ML-runtime.} Сервисы vLLM, TGI, Triton Inference Server и аналогичные не должны требовать модификации: они должны видеть обычный файловый путь в файловой системе и работать со стандартными POSIX-операциями чтения. Переход на новую систему не должен нарушать совместимость с существующими runtime-окружениями.

\textbf{ФТ-4. Поддержка больших файлов весов.} Система должна корректно работать с отдельными файлами весов размером от сотен мегабайт до нескольких сотен гигабайт, включая формат safetensors и бинарные форматы PyTorch. Обязательна поддержка произвольного доступа по смещению (range reads), характерного для ленивой загрузки.

\textbf{ФТ-5. Горизонтальное масштабирование.} Производительность системы не должна деградировать при добавлении новых worker-нод и увеличении числа одновременно запущенных реплик инференса. Добавление узлов должно линейно увеличивать суммарный объём локального кэша, доступного всему кластеру.

\textbf{ФТ-6. Локальность данных и приоритет источников.} При обслуживании запроса на данные должен соблюдаться следующий приоритет источников: первый --- локальный диск той же ноды; второй --- диски соседних нод кластера по внутренней сети; третий --- внешний репозиторий модели. Такая иерархия минимизирует сетевой трафик и задержки.

\subsubsection{Нефункциональные требования}

\textbf{НТ-1. Отказоустойчивость при падении узла.} Выход одной или нескольких worker-нод из строя не должен делать модель полностью недоступной, при условии что её данные присутствуют на других нодах или доступны через fallback на внешний репозиторий. Система должна корректно обнаруживать недоступность узла и переключаться на альтернативный источник.

\textbf{НТ-2. Устранение единой точки отказа (SPOF).} Архитектура не должна содержать единственного центрального компонента, недоступность которого делает систему нефункциональной. Это требование непосредственно исключает использование одного NFS-сервера в качестве единственного хранилища весов.

\textbf{НТ-3. Контролируемые накладные расходы.} Объём метаданных и управляющего трафика (регистрация чанков, heartbeat-сигналы, запросы к реестру) не должен создавать значимой нагрузки на сеть кластера или перегружать реестр метаданных. Накладные расходы на один запрос данных должны составлять пренебрежимо малую долю по сравнению с полезной нагрузкой.

\textbf{НТ-4. Совместимость с Kubernetes.} Все компоненты системы должны развёртываться как стандартные ресурсы Kubernetes --- DaemonSet для агентов на каждой ноде, Deployment для управляющих компонентов, стандартные PVC/PV для постоянного хранения --- без модификации ядра операционной системы, kubelet или компонентов control plane.

Совокупность сформулированных требований определяет целевые свойства архитектуры, которая должна обеспечивать приоритет локальных данных, прозрачность для runtime-сред и отказоустойчивость без центральной точки отказа.

% ============================================================
\subsection{Выводы по разделу}
% ============================================================

Проведённый анализ позволяет сформулировать следующие выводы.

Рост размеров ML-моделей за период с 2012 по 2024~год на три-четыре порядка превратил холодный старт из второстепенной инфраструктурной задержки в архитектурный приоритет. Компоненты $T_{\mathrm{download}}$ и $T_{\mathrm{load}}$ из формулы~(\ref{eq:cold_start}) доминируют в суммарном времени инициализации и напрямую определяются наличием локально кэшированных весов.

Применение закона Литтла~(\ref{eq:little}) показывает, что линейное снижение $W_{\mathrm{cold}}$ пропорционально снижает число ожидающих запросов $L_{\mathrm{wait}}$, что критически важно при burst-масштабировании. Формула~(\ref{eq:pcold}) демонстрирует, что для редко запрашиваемых моделей вероятность холодного старта стремится к единице вне зависимости от политики keep-alive, что оправдывает применение архитектурных механизмов ускорения.

Классическая схема с \texttt{emptyDir} не переиспользует уже загруженные веса между репликами и нодами, что при масштабировании приводит к $N$-кратному увеличению нагрузки на внешний репозиторий. Сформулированные функциональные и нефункциональные требования определяют необходимые свойства системы: локальность данных, прозрачность для ML-runtime, горизонтальная масштабируемость, отказоустойчивость без SPOF.

Центральный вывод раздела состоит в том, что ускорять необходимо не вычислительную фазу инференса, а подготовительную: доставку, кэширование и представление весов модели приложению. Это требует специализированной архитектуры хранения данных, ориентированной на локальность, P2P-передачу между нодами и прозрачный файловый интерфейс для runtime-окружений.
